\chapter{Summary of the pipeline and of the changes}
\label{sec:summary}

\section*{Reductions compared}
Five reductions of the same exposures are compared in this report (Table~\ref{tab:versions_names}): two are
\banzai{} as distributed, the other three successive improved versions of it. These names are used throughout.

\begin{table}[!htbp]
\centering
\caption{The reductions compared in this report.}
\label{tab:versions_names}
\small
\begin{tabular}{@{}lp{11.3cm}@{}}
\toprule
Name & What it is \\
\midrule
\banzai{}~1.0.1 (archive) & The products in the LCO archive. \\
\banzai{}~1.1.0 (stock) & The current release, run here from the raw frames with its default settings, except
that the object wavelength zero point of this work (Section~\ref{sec:wl_zp}) was already applied; without it the
[O\,{\sc iii}] position of Fairall~9 scatters by 6.6\,\AA{} (red order) and 3.1\,\AA{} (blue) between exposures,
instead of 0.6 and 0.4\,\AA. Its fringe, telluric and flux results are those of 1.1.0. \\
Improved v1 & The corrections of Chapters~\ref{sec:traces}--\ref{sec:flux} applied to the archive products
(re-extraction, harmonic fringe model, arc refit, telluric zero point, template$\times$nightly flux
calibration). \\
Improved v2 & The first improved \banzai{} run from raw: harmonic fringe stage (bands masked, onset 6500\,\AA, nodes
every 400\,\AA), combined trace model, telluric correction restricted to the bands and taken from the nine nearest
standards; single-standard flux calibration. \\
\textbf{Improved v3 (final)} & Our final version for now, and the default configuration of the branch: v2 plus the
joint fringe + telluric fit, the residual fringe template of the slit, wider telluric bands, the O$_2$ A band flagged
in science spectra, the ensemble flux calibration, the unused arc lines added to the arc fit, an object-trace template
for each site, and Feige\,67 and GD\,153 as flux standards. Verified at both FLOYDS sites (Chapter~\ref{sec:sites}). \\
\bottomrule
\end{tabular}
\end{table}

\section*{Overview}

This report examines every step of the FLOYDS reduction that matters for reverberation mapping (RM), shows for each
step what the archive pipeline \banzai{}~1.0.1 does, where it fails (with examples from the Fairall~9 campaign), what we
changed, and what the change achieves. The corrections are target-independent: no information about the science
target is used. They were developed on the public \texttt{-1d}/\texttt{-2d} archive files
(Chapters~\ref{sec:traces}--\ref{sec:flux}); the fringe and trace corrections now run inside \banzai{} from the
raw frames, in one pass and with fewer steps than \banzai{}~1.1.0 (stock) (Chapter~\ref{sec:raw}), as do the
telluric correction, the ensemble flux calibration with the grey level from the clearest standards and the added arc
lines. Improved v3 is our final version for now; Chapter~\ref{sec:sites} verifies it on a second campaign taken with
the other FLOYDS (OGG).
Fig.~\ref{fig:flow} shows the chain and Table~\ref{tab:stages} summarises each step. The quality metrics of the
archive products and of the recalibrated spectra are collected in Table~\ref{tab:summary}.



\begin{figure}[!htbp]
\centering
\begin{tikzpicture}[
  node distance=5mm and 7mm,
  box/.style={draw=avaInkBlue, rounded corners=2pt, thick, align=left, font=\small, text width=6.1cm,
              inner sep=4pt, fill=white},
  new/.style={box, fill=blue!6},
  inp/.style={draw=black!45, dashed, rounded corners=2pt, align=left, font=\footnotesize, text width=3.6cm,
              inner sep=3pt, fill=black!3},
  arr/.style={-{Stealth[length=2.2mm]}, thick, avaInkBlue},
  side/.style={-{Stealth[length=1.8mm]}, black!50}]
\node[box] (raw) {\textbf{Raw frames} (science, standards, arcs, lamp flats, sky flats)};
\node[new, below=of raw] (trace) {\textbf{1. Order and object traces} (Ch.~\ref{sec:traces})\\
  order position per block from the same-block lamp flat; object trace checked against the data};
\node[new, below=of trace] (rex) {\textbf{2. Extraction without 2-d fringe division} (Ch.~\ref{sec:fringe})\\
  no flat division; \banzai{} sky model and optimal extraction};
\node[new, below=of rex] (fr1) {\textbf{3. Harmonic fringe fit} (Ch.~\ref{sec:fringe})\\
  phase from lamp-flat master, fitted per spectrum jointly with a telluric template; residual template per slit};
\node[new, below=of fr1] (wl) {\textbf{4. Wavelengths} (Ch.~\ref{sec:wavelength})\\
  arc fit with the added lines; object zero point from the telluric bands};
\node[new, below=of wl] (flux) {\textbf{5. Flux and telluric calibration} (Ch.~\ref{sec:flux})\\
  grey-normalised median of the nine nearest standards (template $\times$ nightly for coarse references)};
\node[new, below=of flux] (fr2) {\textbf{6. Red end of the fringe fit} (Ch.~\ref{sec:fringe})\\
  beyond 8850\,\AA: fit on the telluric-corrected spectrum (one stage in the pipeline)};
\node[box, below=of fr2] (out) {\textbf{Calibrated spectra per order} (red and blue kept separate)};
\foreach \a/\b in {raw/trace, trace/rex, rex/fr1, fr1/wl, wl/flux, flux/fr2, fr2/out} \draw[arr] (\a) -- (\b);
\node[inp, right=of trace] (i1) {sky flats (order shape, one per slit and epoch); same-block lamp flat (position)};
\node[inp, right=of fr1] (i2) {lamp-flat phase master (median over 1--2 months, per slit): fringe \emph{phase} only};
\node[inp, right=of wl] (i3) {raw arcs (Hg\,4916, Ar\,6416/6677/6753 added)};
\node[inp, right=of flux] (i4) {standards (nine nearest in time, any slit); reference fluxes};
\foreach \a/\b in {i1/trace, i2/fr1, i3/wl, i4/flux} \draw[side] (\a) -- (\b);
\node[inp, left=of rex, text width=3.4cm] (q) {QA at every step: narrow lines, sky and telluric zero
  points, order overlap, per-spectrum residuals (Ch.~\ref{sec:diagnostics})};
\draw[side] (q) -- (rex);
\end{tikzpicture}
\caption{The recalibration chain. Shaded boxes are steps that differ from \banzai{}~1.0.1; dashed boxes are the
calibration inputs. The steps act on each order separately.}
\label{fig:flow}
\end{figure}

\begin{table}[!htbp]
\small
\caption{What each step does in \banzai{}~1.0.1, what goes wrong, what we do instead, and the effect.}
\label{tab:stages}
\begin{tabularx}{\textwidth}{@{}>{\raggedright\arraybackslash}p{2.0cm}*{3}{>{\raggedright\arraybackslash}X}@{}}
\toprule
Step & \banzai{}~1.0.1 & Problem & This work \\
\midrule
Order trace &
Order shape from a sky flat (hard-coded list, one per slit and epoch); per-frame vertical shift fitted on the science
frame (red order only). &
None for Fairall~9 (the shift agrees with the lamp flat to 0.14\,px), but the shift is fitted on the science light,
which fails for faint targets. &
Shift from the same-block lamp flat (signal-to-noise ratio in the thousands); the blue order is only traceable on lamp
flats at $x>1300$. \\
Object trace &
Polynomial in wavelength (degree 7 in 1.0.1, 5 with a curvature penalty in 1.1.0) through per-chunk centroids. &
Offsets of 0.8--1.4\,px at the faint red end (9250--9750\,\AA), 1--2 per cent flux loss. &
Physical model: offset + atmospheric differential refraction along the slit + instrumental shape + penalised
correction (default in the pipeline; a template per site). Equal for bright targets; holds at 0.1--0.3\,px where
the polynomial fails ($\gtrsim30\times$ fainter than Fairall~9). \\
Fringing &
Divides each frame by a whole 8-month-old lamp flat, shifted along the slit by up to hundreds of pixels. &
Invalid flat edges destroy data in 13/50 frames; slit illumination biases the flux by up to $\pm10$ per cent; lamp
fringes are smoothed by the 6\,arcsec slit and do not match the object's. &
No 2-d division: a two-harmonic fringe model with the phase from a lamp-flat master is fitted to each extracted
spectrum, jointly with a telluric template, plus a static residual template per slit (for archive products the
division is undone first, Appendix~\ref{sec:app_v1}). Red fringing \FrOld{}$\rightarrow$\FrNewOpt{} per cent
(2.7 in the pipeline from raw). \\
Sky and extraction &
Kelson-type 2-d background; optimal extraction. &
Sky wrong where the flat was invalid; error bars not propagated through the flat. &
Robust per-bin sky with adaptive object mask; optimal and fixed-aperture extractions with propagated errors. \\
Wavelengths &
Arc fit without lines between Hg\,4358 and 5461\,\AA{} (blue) or blueward of 5461\,\AA{} (red); no object zero
point. &
H$\beta$/[O\,{\sc iii}] interpolated or extrapolated ($-8$\,\AA{} in the red at 4916\,\AA); the object's position
in the wide slit shifts its wavelengths by up to 3\,\AA. &
The unused lines added to the arc fit; zero point from the telluric bands in the object's own light, transferred
to the blue order. Narrow-line scatter 160--340$\rightarrow$20--40\,\kms; orders agree to 0.3\,\AA{} (COJ). \\
Flux calibration &
Standard chosen by block/proposal, not slit width; EG\,21 missing from the catalogue; fine-knot spline sensitivity. &
Spikes and fringes in the sensitivity; a single standard's clouds and slit losses transferred to the science. &
Median of the grey-normalised sensitivities of the nine nearest standards, grey level from the clearest of them
(in the pipeline); template $\times$ nightly correction for coarse references; Feige\,67 and GD\,153 added to the
catalogue. Red/blue [O\,{\sc iii}] scatter \FRBOld{}$\rightarrow$\FRBBox{} per cent (6.9 in v3 from raw). \\
Telluric &
Measured on the standard, scaled with airmass. &
Not matched to the seeing-set line-spread function of the 6\,arcsec spectra; red end beyond 8900\,\AA{}
extrapolated. &
Transmission applied only inside the bands, median of the nine nearest standards; O$_2$ A band flagged. Open: a
model transmission matched to each spectrum's resolution and water column (Chapter~\ref{sec:recommendations}). \\
\bottomrule
\end{tabularx}
\end{table}

Table~\ref{tab:summary} collects the main numbers of the diagnosis on the archive data
(Chapters~\ref{sec:traces}--\ref{sec:flux}). All its comparisons use the 50 Fairall~9 spectra and the 18 EG\,21
spectra. ``\banzai{}'' denotes the archive products of \banzai{}~1.0.1. ``Improved v1'' denotes our full
recalibration of these archive products (Appendix~\ref{sec:app_v1}): re-extraction with the new fringe correction and
sky model, arc refit, telluric zero point, and flux calibration. The flux calibration multiplies a sensitivity template
by a nightly correction, using the median of the EG\,21 sensitivities, with the grey level set by the clearest half of
the standards.


\begin{table}[!htbp]
\centering
\caption{Summary of the quality metrics. Scatters are robust (1.4826 times the median absolute deviation) over the
50 Fairall~9 spectra unless stated otherwise. Small-scale scatter: rms over the 50 spectra of the ratio of each
(normalised) spectrum to the median spectrum after removing a 200\,\AA{} running median. Fringe metric: rms of the
ratio of 25\,\AA{} to 200\,\AA{} running medians in the telluric-free windows 7750--8050 and 8450--8900\,\AA.
Optimal and boxcar: the optimal and the fixed-aperture extraction of improved v1 (Appendix~\ref{sec:app_v1}).}
\label{tab:summary}
\begin{tabular}{@{}>{\raggedright\arraybackslash}p{6.9cm}ccc@{}}
\toprule
Metric & \banzai{}~1.0.1 & improved v1 & improved v1 \\
 & (archive) & (optimal) & (boxcar) \\
\midrule
\multicolumn{4}{l}{\textit{Wavelengths (narrow lines at $z=0.047016$; median offset and scatter, \kms)}}\\
{[O\,{\sc iii}]\,5007, blue order} & \OIIIbOld & \multicolumn{2}{c}{\OIIIbNew} \\
{[O\,{\sc iii}]\,5007, red order} & \OIIIrOld & \multicolumn{2}{c}{\OIIIrNew} \\
{[O\,{\sc iii}]\,4959, blue order} & \OIIIfourOld & \multicolumn{2}{c}{\OIIIfourNew} \\
{[Ne\,{\sc iii}]\,3869, blue order} & \NeIIIOld & \multicolumn{2}{c}{\NeIIINew} \\
{[O\,{\sc i}]\,6300, red order} & \OIOld & \multicolumn{2}{c}{\OINew} \\
Ca\,{\sc ii}\,K (host), blue order & \CaKOld & \multicolumn{2}{c}{\CaKNew} \\
\midrule
\multicolumn{4}{l}{\textit{Red-order fringing and artefacts (units in each row)}}\\
Fringe metric (median of 50, per cent) & \FrOld & \FrNewOpt & \FrNewBox \\
Fringe in the median of 50 spectra, 7300--9500\,\AA{} (per cent) & \FrMedOld & \FrMedNewOpt & \FrMedNewBox \\
$>8\sigma$ spikes at 4600--5800\,\AA{} (number per spectrum, median / maximum) & \SpOld & \SpNewOpt & \SpNewBox \\
\midrule
\multicolumn{4}{l}{\textit{Small-scale scatter between spectra (per cent)}}\\
blue 3500--4000\,\AA & \SSbaOld & \SSbaOpt & \SSbaBox \\
blue 4000--4600\,\AA & \SSbbOld & \SSbbOpt & \SSbbBox \\
blue 4600--5200\,\AA & \SSbcOld & \SSbcOpt & \SSbcBox \\
blue 5200--5700\,\AA & \SSbdOld & \SSbdOpt & \SSbdBox \\
red 5000--6000\,\AA & \SSraOld & \SSraOpt & \SSraBox \\
red 6000--6800\,\AA & \SSrbOld & \SSrbOpt & \SSrbBox \\
red 7000--7500\,\AA & \SSrcOld & \SSrcOpt & \SSrcBox \\
red 7700--8500\,\AA & \SSrdOld & \SSrdOpt & \SSrdBox \\
red 8500--9500\,\AA & \SSreOld & \SSreOpt & \SSreBox \\
\midrule
\multicolumn{4}{l}{\textit{Flux consistency (per cent)}}\\
{[O\,{\sc iii}]\,5007 flux scatter, blue order} & \FOIIIbOld & \FOIIIbOpt & \FOIIIbBox \\
{[O\,{\sc iii}]\,5007 flux scatter, red order} & \FOIIIrOld & \FOIIIrOpt & \FOIIIrBox \\
red/blue [O\,{\sc iii}] flux ratio scatter & \FRBOld & \FRBOpt & \FRBBox \\
EG\,21 calibrated shape vs CTIO reference, blue, rms (18 nights) & \CLbOld & \multicolumn{2}{c}{\CLbNew} \\
EG\,21 calibrated shape vs CTIO reference, red, rms (18 nights) & \CLrOld & \multicolumn{2}{c}{\CLrNew} \\
EG\,21 red fringe-scale rms (median of 18) & \EGfrOld & \multicolumn{2}{c}{\EGfrNew} \\
\bottomrule
\end{tabular}
\end{table}

